% Ad Familiares 14.4 --- headnote
% ad-familiares-14-04, 29 April 58 BC, written at Brundisium

\textit{Cicero to his wife Terentia, his daughter Tullia, and
his son Marcus, written from Brundisium on the eve of his
departure across the Adriatic into exile, late April 58 BC
(the body gives \textit{a. d. ii K. Mai.} ---  30 April --- as
the date of departure, and \textit{Pr. K. Mai.} for the
sealing of the letter; the Perseus dateline-header reads
\textit{Mart.} and is generally held to be a scribal error
for \textit{Mai.}). The first surviving letter of the
fourteen letters to Terentia: ten of them from this exile,
the rest from the long uneasiness of the post-exile years and
the first months of the civil war. The voice is the consular's
voice as it has not been heard before in the corpus: undone,
weeping, unable to write or read.}

\textit{The substance of the letter: thirteen days have been
spent at Brundisium with M. Laenius Flaccus, who took Cicero
in despite the penalty of the lex Clodia (any harbouring of
the proscribed counted as a capital offence). Cicero is
sailing tomorrow for Cyzicus through Macedonia, the
Thessalonican destination not yet decided; the wind and the
sailors will not wait. He cannot bear to ask Terentia, whom
he calls \textit{aegram et corpore et animo confectam}, to
come to him; nor can he bear to think of being without her;
the matter must go as it goes. The poignant centre is the
two children: Tullia (now married to C. Calpurnius Piso
Frugi, who features in \textsection 4 as ``Piso'') ---
``the poor little thing's marriage and her good name must be
guarded''; and the boy Cicero, ``always in my bosom and
embrace.'' \textsection 5 closes on the great line ---
\textit{viximus, floruimus; non vitium nostrum sed virtus
nostra nos adflixit} (``we have lived, we have flourished;
not our vice but our virtue has cast us down'') --- and on
the bare admission: ``I, who am putting heart into you,
cannot put it into myself.''}
